
\medskip
\phantomsection\label{ex:sga3-VIB-5.8.3}%
\noindent\textbf{Example 5.8.3.---}
Fix an algebraically closed field \(k\). Recall first that every
``locally Boolean'' topological space \(X\), that is, every space
having a basis of compact open subsets, becomes a \(k\)-scheme
when equipped with the sheaf of locally constant \(k\)-valued
functions (cf.\ [DG70], I, \S1, \SGARefLink{cor:sga3-VIB-2.12}{2.12}). Such an \(X\) will be
called a \emph{locally Boolean \(k\)-scheme}.

We also note that every topological space \(E\) having a basis of
separated open subsets---and likewise every scheme \(X\)---has a
dense separated open subset. Indeed, every increasing union of
separated open subsets is separated (for a scheme \(X\), this
follows from \hyperref[lem:sga3-VIB-5.6.3]{5.6.3}), so there is a maximal such open subset \(U\).
It must be dense: otherwise there would be a nonempty open subset
\(V\) with \(U\cap V=\varnothing\); the open set \(V\) would
contain a nonempty separated open subset \(W\), and
\[
  U\cup W
\]
would still be separated, contradicting the maximality of \(U\).

Now let \(C\) be the Cantor space, viewed as the underlying space
of the group
\[
  (\mathbb Z/2\mathbb Z)^{\mathbb N}
\]
with its product topology. For every \(p\in C\), let \(C(p)\) be
another copy of \(C\), and let \(X\) be obtained by gluing each
\(C(p)\) to \(C\) along \(C-\{p\}\). Then \(X\) is a
nonseparated locally Boolean \(k\)-scheme, and \(C\) is a dense
open subset of \(X\).

Let \(G\) be the group of automorphisms of the \(k\)-scheme \(X\),
that is, of homeomorphisms of \(X\). Then \(G\) acts transitively
on \(X\). Indeed, \(X\) is the union of \(C\) and, for each point
\(p\in C\), a second point \(p'\). This second point is
characterized as the unique \(x\in X-\{p\}\) for which
\[
  (C-\{p\})\cup\{x\}
\]
is separated. Consequently, every automorphism \(\phi\) of \(C\)
extends to an automorphism \(\phi_X\) of \(X\) satisfying
\[
  \phi_X(p')=\phi_X(p)'
\]
for every \(p\). Moreover, the map
\[
  \theta_p:X\longrightarrow X
\]
that exchanges \(p\) and \(p'\) and fixes every other point is an
automorphism of \(X\). Finally, the automorphism group of \(C\)
acts transitively on \(C\): using the group structure of \(C\),
it is enough to consider translations.

Thus the discrete \(k\)-group \(G\), which is locally of finite
type, acts on the \(k\)-scheme \(X\) in such a way that
\[
  G\times X\longrightarrow X\times X,
  \qquad (g,x)\longmapsto(gx,x),
\]
is surjective, but \(X\) is not separated, even though \(C\) is a
dense separated open subset.
